% Source: https://artofproblemsolving.com/community/c3h2435540_amc10_mock_problems
% Stable identity: aops-topic-2435540-q06
Let $O$ be the centre of the large gear and $Q$ the centre of the small gear, so $OQ=103.5+34.5=138$ and $Q$ moves on a circle about $O$.

\begin{asy}
usepackage("amsmath");
size(5cm);
defaultpen(fontsize(10pt));
real R=3, r=1;
pair O=(0,0), Q=(R+r)*dir(50), T=R*dir(50);
draw(circle(O,R),linewidth(1)); draw(circle(Q,r),blue+linewidth(1));
draw(circle(O,R+r),dashed+gray);
draw(O--(R+r+0.4,0),dotted); draw(O--Q,dashed);
dot(O); dot(Q); dot(T);
label("$O$",O,SE); label("$Q$",Q,NE,blue); label("$T$",T+0.12*dir(192),dir(192));
draw(O--R*dir(210),Arrow(4)); label("$103.5$",0.5R*dir(210),dir(300));
draw(Q--Q+r*dir(-30),blue,Arrow(4)); label("$34.5$",Q+r*dir(-30),E,blue);
draw(arc(O,1.0,0,50),Arrow(4)); label("$\theta$",1.25*dir(25));
label("path of $Q$",(R+r)*dir(-45),SE,gray);
\end{asy}

Suppose the line $OQ$ has turned through an angle $\theta$ (in radians). The contact point $T$ has moved along an arc of length $103.5\,\theta$ of the large gear, and since there is no slipping, the same length of the small gear's rim has rolled over it. So relative to the line $OQ$ the small gear has turned through $\dfrac{103.5\,\theta}{34.5}=3\theta$, and since the line $OQ$ itself has turned through $\theta$ in the same direction, the small gear has rotated through
\[3\theta+\theta=4\theta.\]
Three complete rotations mean $4\theta=3\cdot2\pi$, i.e. $\theta=\dfrac34\cdot2\pi$. So the small gear has gone $\boxed{\dfrac{3}{4}}$ of a revolution around the larger gear. (In one full trip around it rotates $\dfrac{103.5+34.5}{34.5}=4$ times.)
